\chapter{Inference in Regression}\label{ch:b2:regression-inference}

A regression coefficient is an estimate; its standard error says how far it might be from what it estimates. The standard
errors printed by default rest on assumptions about the errors that real data often break: that they all have the same
variance, and that they are independent of each other. This chapter states those assumptions, shows why least squares is
the best linear estimator when they hold, and builds the standard errors that remain valid when they do not. The chart
shows how much the answer can change.

\begin{openingchart}{Heteroskedasticity-robust standard error divided by the classical standard error, for each
  coefficient of a regression of outpatient visits to a doctor on insurance-plan and health variables; 20,190
  person-years.}{RAND Health Insurance Experiment, statsmodels randhie dataset; our computation}
\begin{tikzpicture}
\begin{axis}[omrecord, width=0.72\linewidth, height=0.4\linewidth, ybar, bar width=10pt, xmin=-0.6, xmax=9.6,
    ymin=0, ymax=1.8, xtick={0,...,9}, xticklabels from table={figdata/bachelor-2/31-regression-inference/out/se.csv}{label},
    xticklabel style={rotate=35, anchor=north east, font=\scriptsize}, ylabel={Robust / classical},
    nodes near coords, nodes near coords style={font=\scriptsize, fill=white, inner sep=1pt, yshift=1pt,
    /pgf/number format/.cd, fixed, precision=2}]
  \draw[guide] (axis cs:-0.6,1) -- (axis cs:9.6,1);
  \addplot[fill=omDef!70, draw=omDef] table[x expr=\coordindex, y=ratio]
    {figdata/bachelor-2/31-regression-inference/out/se.csv};
\end{axis}
\end{tikzpicture}
\end{openingchart}

% figdata: se (ratio hlthp 1.55, disea 1.29, physlm 1.25, hlthf 1.18, lncoins 0.93, hlthg 0.92)
For most coefficients the two standard errors are within 10 per cent of each other. For four they are not: the robust
standard error is 1.55 times the classical one for poor self-rated health, 1.29 times for the number of chronic diseases,
1.25 times for a physical limitation and 1.18 times for fair health. The classical formula assumes that every person's visits are equally
unpredictable; the sick are much less predictable than the healthy.

\begin{recall}
From the Year 1 volume: estimators, unbiased and consistent estimators, standard errors, confidence intervals, the $t$ and
$F$ distributions, tests and $p$-values, and the central limit theorem. From \cref{ch:b2:multiple-regression}: the
matrix form, the \omterm{def:b2:multiple-regression:normal-equations}{normal equations} and the \omterm{def:b2:multiple-regression:fwl}{Frisch--Waugh--Lovell theorem}.
\end{recall}

\section{The Gauss--Markov theorem}\label{sec:b2:regression-inference:gm}

\begin{definition}[Gauss--Markov assumptions]\label{def:b2:regression-inference:gm-assumptions}
The \emph{Gauss--Markov assumptions}\index{Gauss--Markov assumptions} for $y = X\beta + u$ are: $X$ has full column rank;
$E[u \mid X] = 0$; and $\operatorname{Var}(u \mid X) = \sigma^2I$, so that the errors are uncorrelated with each other and
all have the same variance.
\end{definition}

\begin{definition}[Homoskedasticity]\label{def:b2:regression-inference:homoskedasticity}
Errors are \emph{homoskedastic}\index{homoskedasticity} if their variance given the regressors is the same for every
observation, $\operatorname{Var}(u_i \mid X) = \sigma^2$; they are \emph{heteroskedastic}\index{heteroskedasticity} if it
varies across observations.
\end{definition}

Under the first two assumptions $\hat\beta = (X'X)^{-1}X'y = \beta + (X'X)^{-1}X'u$ is unbiased. Under the third,
$\operatorname{Var}(\hat\beta \mid X) = (X'X)^{-1}X'(\sigma^2I)X(X'X)^{-1} = \sigma^2(X'X)^{-1}$, and $\sigma^2$ is
estimated by $s^2 = e'e/(n - k)$. These are the classical standard errors.

\begin{definition}[Best linear unbiased estimator]\label{def:b2:regression-inference:blue}
An estimator of $\beta$ is the \emph{best linear unbiased estimator}\index{best linear unbiased estimator} if it is linear
in $y$, unbiased, and has a variance matrix no larger, in the sense of positive semi-definite matrices, than that of any
other linear unbiased estimator.
\end{definition}

\begin{theorem}[Gauss--Markov]\label{thm:b2:regression-inference:gauss-markov}
Under the \omterm{def:b2:regression-inference:gm-assumptions}{Gauss--Markov assumptions}, the least-squares estimator is the \omterm{def:b2:regression-inference:blue}{best linear unbiased estimator} of $\beta$.
\end{theorem}

\begin{proof}
Let $\tilde\beta = Cy$ be linear, and write $C = (X'X)^{-1}X' + D$. Then $E[\tilde\beta \mid X] = \beta + DX\beta$, so
unbiasedness for every $\beta$ requires $DX = 0$. Then $\tilde\beta - \beta = Cu$, and
$\operatorname{Var}(\tilde\beta \mid X) = \sigma^2CC' = \sigma^2\bigl((X'X)^{-1} + (X'X)^{-1}X'D' + DX(X'X)^{-1} +
DD'\bigr) = \sigma^2(X'X)^{-1} + \sigma^2DD'$, since $DX = 0$. The difference from the least-squares variance is
$\sigma^2DD'$, which is positive semi-definite.
\end{proof}

The listing compares least squares with another linear unbiased estimator of a slope: the difference between the mean
outcomes of the upper and lower halves of the regressor, divided by the difference of their mean regressors.

\omcode{code/bachelor-2/31-regression-inference/listings/gauss_markov_sim.py}{1}{29}{Least squares and the grouping
estimator of a slope, over 20,000 simulated samples.}\label{lst:b2:regression-inference:gm}

\omoutput{code/bachelor-2/31-regression-inference/outputs/gauss_markov_sim.txt}{Both are centred on the true slope 0.5;
least squares has the smaller variance, as the theorem says.}\label{out:b2:regression-inference:gm}

% code: listings/gauss_markov_sim.py
The variances are 0.01145 and 0.01508 in the simulation, against 0.01141 and 0.01521 by the formulas. The theorem says
nothing about estimators that are biased or non-linear, and nothing at all once the errors are \omterm{def:b2:regression-inference:homoskedasticity}{heteroskedastic}.

\section{Normal errors: $t$ and $F$}\label{sec:b2:regression-inference:tf}

If, in addition, the errors are normal given $X$, then $(\hat\beta_j - \beta_j)/\operatorname{se}(\hat\beta_j)$ has a $t$
distribution with $n - k$ degrees of freedom, and confidence intervals and tests on one coefficient follow as in the Year 1
volume. Several restrictions at once need another statistic.

\begin{definition}[Joint hypothesis]\label{def:b2:regression-inference:joint-hypothesis}
A \emph{joint hypothesis}\index{joint hypothesis} is a set of $q \geq 2$ restrictions on the coefficients tested together,
such as $\beta_2 = \beta_3 = \beta_4 = 0$.
\end{definition}

\begin{definition}[F statistic]\label{def:b2:regression-inference:f-statistic}
The \emph{F statistic}\index{F statistic} for $q$ linear restrictions is $F = \dfrac{(\mathit{SSR}_r -
\mathit{SSR}_u)/q}{\mathit{SSR}_u/(n - k)}$, where $\mathit{SSR}_r$ and $\mathit{SSR}_u$ are the sums of squared residuals of the
regressions with and without the restrictions. Under the classical assumptions with normal errors it has an $F(q, n - k)$
distribution when the restrictions hold.
\end{definition}

The $F$ statistic asks whether imposing the restrictions raises the sum of squared residuals by more than chance would. The
listing applies it to the RAND Health Insurance Experiment, in which families were assigned at random to insurance plans
that made them pay different shares of their medical bills.

\omcode{code/bachelor-2/31-regression-inference/listings/t_and_f.py}{1}{25}{Doctor visits on plan and health variables;
the joint test of the three health dummies, by hand and with statsmodels.}\label{lst:b2:regression-inference:tf}

\omoutput{code/bachelor-2/31-regression-inference/outputs/t_and_f.txt}{Coefficient, classical standard error and $t$
statistic for each term; the $F$ statistic of the joint test.}\label{out:b2:regression-inference:tf}

% code: listings/t_and_f.py; ledger: b2:randhie-data, b2:randhie2006
Each unit of log coinsurance goes with 0.17 fewer visits, holding the other plan variables and health fixed ($t = -8.41$).
The health dummies, measured against excellent health, have $t$ statistics of $-0.73$, $1.81$ and $5.53$; jointly they
raise the sum of squared residuals by 667.9 and give $F = 11.78$, far above the 5 per cent critical value of 2.605. Both
calculations used the classical variance. The next section asks whether it was right to.

\section{Unequal error variances}\label{sec:b2:regression-inference:robust}

\begin{omchart}{Standard deviation of the residuals in twenty equal groups of person-years sorted by fitted number of
  visits, against the mean fitted value of each group.}{RAND Health Insurance Experiment, statsmodels randhie dataset; our
  computation}
\begin{tikzpicture}
\begin{axis}[omrecord, width=0.6\linewidth, height=0.38\linewidth, xmin=0.5, xmax=6.5, ymin=0, ymax=9,
    xlabel={Mean fitted visits}, ylabel={S.d. of residuals}]
  \addplot[omDef, thick, mark=*, mark size=1.3pt] table[x=fitted, y=sd]
    {figdata/bachelor-2/31-regression-inference/out/spread.csv};
\end{axis}
\end{tikzpicture}
\end{omchart}

% figdata: spread (first bin sd 2.443 at fitted 0.95; last bin 8.197 at fitted 6.135)
Among the person-years with the fewest predicted visits the residuals have a standard deviation of 2.44; among those with
the most, 8.20. The errors are \omterm{def:b2:regression-inference:homoskedasticity}{heteroskedastic}, and the classical variance formula does not apply. Least squares is still
unbiased, but its variance is now
\begin{equation}\label{eq:b2:regression-inference:sandwich}
\operatorname{Var}(\hat\beta \mid X) = (X'X)^{-1}\Bigl(\sum_i\sigma_i^2x_ix_i'\Bigr)(X'X)^{-1},
\end{equation}
with $x_i'$ the $i$th row of $X$ and $\sigma_i^2 = \operatorname{Var}(u_i \mid X)$ (Exercise 11). Its shape gives it the
name of a sandwich.

\begin{definition}[Heteroskedasticity-robust standard error]\label{def:b2:regression-inference:robust-se}
A \emph{heteroskedasticity-robust standard error}\index{heteroskedasticity-robust standard error} is the square root of a
diagonal element of \eqref{eq:b2:regression-inference:sandwich} with each $\sigma_i^2$ replaced by the squared residual
$e_i^2$ (HC0), often multiplied by $n/(n - k)$ (HC1).
\end{definition}

Each $e_i^2$ is a poor estimate of its $\sigma_i^2$, but the sum in the middle averages over many observations, and in
large samples the robust variance is consistent whatever the form of \omterm{def:b2:regression-inference:homoskedasticity}{heteroskedasticity}.

\omcode{code/bachelor-2/31-regression-inference/listings/robust_se.py}{1}{27}{The sandwich by hand, against
statsmodels.}\label{lst:b2:regression-inference:robust}

\omoutput{code/bachelor-2/31-regression-inference/outputs/robust_se.txt}{Classical, HC0 and HC1 standard errors; with
20,190 observations HC0 and HC1 differ only in the fourth decimal.}\label{out:b2:regression-inference:robust}

What happens to a test when the wrong standard error is used? The next listing draws samples in which the true slope is
zero and the error spread grows with the regressor, and counts how often a nominal 5 per cent test rejects.

\omcode{code/bachelor-2/31-regression-inference/listings/size_sim.py}{1}{18}{The rejection rate of a true null, with
classical and robust standard errors.}\label{lst:b2:regression-inference:size}

\omoutput{code/bachelor-2/31-regression-inference/outputs/size_sim.txt}{The classical test rejects a true null a quarter
of the time; the robust test close to the nominal 5 per cent.}\label{out:b2:regression-inference:size}

% code: listings/size_sim.py
The classical test rejects in 25.4 per cent of samples, five times as often as it claims; the robust test in 5.7 per cent.
With a \omterm{def:b2:regression-inference:joint-hypothesis}{joint hypothesis} the robust version of the $F$ statistic uses the robust variance matrix; on the RAND data the
test of the health dummies falls from 11.78 to 5.40, still above the critical value (Exercise 6).

\section{Clustered errors}\label{sec:b2:regression-inference:cluster}

Observations often come in groups: pupils in a class, workers in a firm, families in a town, the same person in several
years. Anything shared by the group (a teacher, a manager, a local shock) puts a common component in the errors of its
members. If the regressor is also shared within the group, the sample contains much less independent information than its
size suggests.

Suppose there are $G$ groups of $m$ members, the regressor $x$ is the same for every member of a group, and the errors
have variance $\sigma^2$ and correlation $\rho$ between members of the same group. Then (Exercise 12)
\begin{equation}\label{eq:b2:regression-inference:moulton}
\operatorname{Var}(\hat\beta_1 \mid X) = \frac{\sigma^2}{\sum_i(x_i - \bar x)^2}\,\bigl(1 + (m - 1)\rho\bigr).
\end{equation}

\begin{definition}[Moulton factor]\label{def:b2:regression-inference:moulton}
The \emph{Moulton factor}\index{Moulton factor} is the ratio of the true standard error of a coefficient to the classical
one when errors are correlated within groups; for groups of equal size $m$, a regressor constant within groups and an
error correlation $\rho$, it is $\sqrt{1 + (m - 1)\rho}$.
\end{definition}

With groups of 20 and a correlation of only 0.2, the factor is 2.19: classical standard errors are less than half what
they should be.

\begin{definition}[Cluster-robust standard error]\label{def:b2:regression-inference:cluster-se}
A \emph{cluster-robust standard error}\index{cluster-robust standard error} replaces the middle of the sandwich by
$\sum_g\bigl(\sum_{i \in g}e_ix_i\bigr)\bigl(\sum_{i \in g}e_ix_i\bigr)'$, summing within each group first, so that the
errors may be correlated in any way within groups but are independent across them.
\end{definition}

\omcode{code/bachelor-2/31-regression-inference/listings/cluster_sim.py}{1}{26}{Fifty groups of twenty: classical,
robust and \omterm{def:b2:regression-inference:cluster-se}{cluster-robust standard errors} against the true spread of the estimate.}\label{lst:b2:regression-inference:cluster}

\omoutput{code/bachelor-2/31-regression-inference/outputs/cluster_sim.txt}{Only the \omterm{def:b2:regression-inference:cluster-se}{cluster-robust standard error} is close
to the true standard deviation of the estimate.}\label{out:b2:regression-inference:cluster}

% code: listings/cluster_sim.py
Across 500 samples the slope has a standard deviation of 0.0696. The classical standard errors average 0.0323, and the
robust ones 0.0322: robustness to unequal variances does nothing for correlation within groups. The
\omterm{def:b2:regression-inference:cluster-se}{cluster-robust standard error} averages 0.0684, 2.115 times the classical one, close to the \omterm{def:b2:regression-inference:moulton}{Moulton factor} of 2.191.
\omterm{def:b2:regression-inference:cluster-se}{Cluster-robust standard errors} need many clusters: their middle term is an average over groups, and with few groups it is
itself very noisy, and often too small (the weekend problem has three).

The RAND data of this chapter are such a case. The experiment assigned insurance plans to families, not to persons, and
kept each family on its plan for several years, so the plan variables are the same for every member and every year of a
family, and one person's errors are correlated across years. The standard errors of the plan coefficients should then be
clustered by family; the classical and the robust ones printed above are both likely too small. The file used here has
no family or person identifier, so the clustering cannot be done with it.

\section{Large-sample results}\label{sec:b2:regression-inference:asymptotic}

Without normal errors, the $t$ and $F$ distributions are not exact. Large samples give approximations that need neither
normality nor \omterm{def:b2:regression-inference:homoskedasticity}{homoskedasticity}.

\begin{definition}[Asymptotic normality]\label{def:b2:regression-inference:asymptotic-normality}
An estimator $\hat\theta$ is \emph{asymptotically normal}\index{asymptotic normality} if $\sqrt n(\hat\theta - \theta)$
converges in distribution to a normal distribution with mean zero as $n \to \infty$.
\end{definition}

\begin{theorem}[Large-sample distribution of least squares]\label{thm:b2:regression-inference:asymptotic}
If the observations $(y_i, x_i)$ are independent and identically distributed, $E[x_iu_i] = 0$, $Q = E[x_ix_i']$ is
invertible and the relevant moments are finite, then $\hat\beta$ is consistent and $\sqrt n(\hat\beta - \beta)$ converges
in distribution to $N\bigl(0, Q^{-1}E[u_i^2x_ix_i']Q^{-1}\bigr)$.
\end{theorem}
\admitted

The variance in the theorem is the sandwich again; its sample version is the robust variance matrix. So $t$ statistics
built with robust standard errors are approximately standard normal in large samples, which is why the robust test of the
listing rejected close to 5 per cent of the time.

\section{The bootstrap}\label{sec:b2:regression-inference:bootstrap}

\begin{definition}[Bootstrap]\label{def:b2:regression-inference:bootstrap}
The \emph{bootstrap}\index{bootstrap} estimates the sampling distribution of a statistic by recomputing it on many samples
drawn with replacement from the data, treating the sample as if it were the population.
\end{definition}

To \omterm{def:b2:regression-inference:bootstrap}{bootstrap} a regression coefficient, draw $n$ observations (rows) with replacement, run the regression, keep the
coefficient, and repeat a thousand times or more; the standard deviation of the kept coefficients is a standard error.
Drawing rows keeps each observation's regressors and error together, so the result is robust to \omterm{def:b2:regression-inference:homoskedasticity}{heteroskedasticity}; with
grouped data whole groups must be drawn. The \omterm{def:b2:regression-inference:bootstrap}{bootstrap} is useful for statistics without a simple variance formula. It does
not repair a bad design: drawing from a sample that has three clusters still gives three clusters.

\begin{model}{The classical linear model}
\assumes The \omterm{def:b2:regression-inference:gm-assumptions}{Gauss--Markov assumptions} and, for exact $t$ and $F$ distributions, normal errors.
\predicts Least squares is the \omterm{def:b2:regression-inference:blue}{best linear unbiased estimator}; classical standard errors and $t$ and $F$ tests have their
stated properties in any sample size.
\tested The spread of residuals against fitted values, and the comparison of classical with robust and \omterm{def:b2:regression-inference:cluster-se}{cluster-robust
standard errors}: on the RAND data the residual spread rises more than threefold across fitted values.
\breaks \omterm{def:b2:regression-inference:homoskedasticity}{Heteroskedasticity} makes classical standard errors wrong in either direction (25.4 per cent rejections in the
simulation); correlation within groups makes them too small by the \omterm{def:b2:regression-inference:moulton}{Moulton factor}; robust and \omterm{def:b2:regression-inference:cluster-se}{cluster-robust standard
errors} repair these in large samples and with many clusters, not with few.
\end{model}

\section{Exercises}\label{sec:b2:regression-inference:exercises}

\begin{exercise}[$\star$]\label{exo:b2:regression-inference:1}
Read the opening chart. For which coefficients is the robust standard error larger than the classical one by more than 20
per cent? For which is it smaller?
\end{exercise}

\begin{exercise}[$\star$]\label{exo:b2:regression-inference:2}
In the regression of the listing \texttt{t\_and\_f.py}, interpret the coefficient on \texttt{physlm} and its $t$
statistic, with classical and with robust standard errors.
\end{exercise}

\begin{exercise}[$\star$]\label{exo:b2:regression-inference:3}
A test at the 5 per cent level rejects a true null in 25.4 per cent of simulated samples. What does that mean for a
researcher who reports a ``significant'' coefficient with classical standard errors in such data?
\end{exercise}

\begin{exercise}[$\star$]\label{exo:b2:regression-inference:4}
Pupils are taught in classes of 25, and the error correlation within a class is 0.1. A regressor is the same for every
pupil in a class. Compute the \omterm{def:b2:regression-inference:moulton}{Moulton factor}, and say how many independent pupils the sample is worth per class.
\end{exercise}

\begin{exercise}[$\star\star$]\label{exo:b2:regression-inference:5}
From the listing's output, compute the $F$ statistic of the health dummies by hand from the two sums of squared residuals.
\end{exercise}

\begin{exercise}[$\star\star$]\label{exo:b2:regression-inference:6}
\textbf{(D)} With the RAND data, test jointly that the three health dummies are zero and, separately, that the four plan
variables are zero, with classical and with robust variance matrices.
\end{exercise}

\begin{exercise}[$\star\star$]\label{exo:b2:regression-inference:7}
In the Gauss--Markov listing, use the bottom and top quarters of $x$ instead of the halves. Predict, then compute, the
variance of the grouping estimator. Does it beat least squares?
\end{exercise}

\begin{exercise}[$\star\star$]\label{exo:b2:regression-inference:8}
Change the size listing so that the errors are \omterm{def:b2:regression-inference:homoskedasticity}{homoskedastic}. Predict, then compute, the two rejection rates.
\end{exercise}

\begin{exercise}[$\star\star$]\label{exo:b2:regression-inference:9}
In the cluster listing, what would the \omterm{def:b2:regression-inference:moulton}{Moulton factor} be if the errors were uncorrelated within groups? If the regressor
varied within groups and were uncorrelated between members of a group?
\end{exercise}

\begin{exercise}[$\star\star\star$]\label{exo:b2:regression-inference:10}
Show that the least-squares estimator of $\beta_1$ in $y_i = \beta_0 + \beta_1x_i + u_i$ is a linear function of the $y_i$
with weights that sum to zero and satisfy $\sum_iw_ix_i = 1$, and show directly that any other set of such weights gives a
larger variance under \omterm{def:b2:regression-inference:homoskedasticity}{homoskedasticity}.
\end{exercise}

\begin{exercise}[$\star\star\star$]\label{exo:b2:regression-inference:11}
Derive \eqref{eq:b2:regression-inference:sandwich} from $\hat\beta - \beta = (X'X)^{-1}X'u$, and show that it reduces to
$\sigma^2(X'X)^{-1}$ under \omterm{def:b2:regression-inference:homoskedasticity}{homoskedasticity}.
\end{exercise}

\begin{exercise}[$\star\star\star$]\label{exo:b2:regression-inference:12}
Derive \eqref{eq:b2:regression-inference:moulton}.
\end{exercise}

\section{Problem: Brenland's census of firms}\label{sec:b2:regression-inference:problem}

\begin{problem}[{Weekend problem --- Brenland's census of firms: a log production function, robust and clustered
intervals, the test of constant returns, and three clusters}]
\label{pb:b2:regression-inference:1}
Brenland's census of 2025 covers 500 firms in three sectors: 51 in farming, 123 in manufacturing and 326 in services. Let
$y$ be log output, $l$ log workers and $k$ log capital, and estimate $y = \beta_0 + \beta_ll + \beta_kk + u$.

\medskip
\noindent\textbf{Part I --- A log production function.}
\begin{enumerate}
  \item What production function does the regression correspond to?
  \item What do $\beta_l$ and $\beta_k$ measure?
  \item Give the estimates and $R^2$.
  \item What does $\beta_l + \beta_k$ measure?
  \item Why might firms with a high error $u$ also choose more workers, and what would that do to $\hat\beta_l$?
\end{enumerate}

\medskip
\noindent\textbf{Part II --- Robust and clustered intervals.}
\begin{enumerate}[resume]
  \item Give the classical standard error of $\hat\beta_k$.
  \item Give the heteroskedasticity-robust one, and a 95 per cent interval.
  \item Give the standard error clustered by sector.
  \item With three clusters, which critical value should the interval use?
  \item Give the clustered 95 per cent interval.
  \item Compare the three standard errors. Which is smallest, and should that reassure?
\end{enumerate}

\medskip
\noindent\textbf{Part III --- The test of constant returns.}
\begin{enumerate}[resume]
  \item Write the hypothesis of constant returns as a restriction.
  \item Write the restricted regression.
  \item Give the two sums of squared residuals.
  \item Compute the $F$ statistic.
  \item Compare it with the 5 per cent critical value of $F(1, 497)$, 3.86.
  \item Give the robust version of the test.
  \item State the verdict.
\end{enumerate}

\medskip
\noindent\textbf{Part IV --- What changes with three clusters.}
\begin{enumerate}[resume]
  \item Why is a cluster-robust variance from three clusters unreliable?
  \item How many degrees of freedom does a $t$ test with three clusters have?
  \item Why can the clustered standard error be smaller than the robust one here?
  \item What could replace sector clustering with so few sectors?
  \item If the census had 300 villages, each with a few firms, would clustering by village help?
  \item Does the \omterm{def:b2:regression-inference:bootstrap}{bootstrap} solve the problem of three clusters?
  \item State the result: the capital elasticity with robust and clustered intervals, and the verdict on constant returns.
\end{enumerate}
\end{problem}
